\begin{remarks}[7.2.10]
\label{IV.2.s7.2.10-fr}
\begin{enumerate}
\item[(i)] De (7.2.7) y (7.2.8) se deduce que todo anillo cociente de
un anillo local de Cohen-Macaulay es estrictamente formalmente
catenario.

Un anillo local noetheriano de dimensión 2 que satisfaga
$(\mathrm S_2)$ es estrictamente formalmente catenario, puesto que es
un anillo de Cohen-Macaulay. Recordemos, en cambio, que existen
anillos locales noetherianos íntegros de dimensión 2 que no son
universalmente catenarios ni, \emph{con mayor razón}, estrictamente
formalmente catenarios (5.6.11).
\item[(ii)] No se sabe si un anillo formalmente catenario (7.1.9) es
estrictamente formalmente catenario; esto se debe a que se desconoce
si, para un anillo local noetheriano íntegro $B$, su completado
$\widehat B$ carece de ideales primos asociados inmersos (6.4.3).
Ignoramos asimismo si, en las condiciones equivalentes \emph{c)}, \emph{d)}, \emph{e)} de
(7.2.5), puede sustituirse la hipótesis de que $A$ es universalmente
catenario por la de que $A$ es catenario; se puede demostrar que así
ocurre cuando $A$ es henseliano (18.9.6).
\end{enumerate}
\end{remarks}
